训练日志上，损失从 \(0.412\) 降到 \(0.409\)。这算降了吗？

问题不在这两个数，而在你手里没有一把尺子。你想知道的其实是变化的速度——它还在往下走，还是已经贴着地面滑行？可``速度''这个词，在只有一串离散数值的地方并没有定义。取相邻两步算差商，得到的数随步长剧烈摇摆；步长取得越小，分子分母同时趋于零，比值反而更不稳定。想问的问题很清楚，能用的算术却答不了。

微积分处理的就是这一类困境：\textbf{你关心的量，恰好落在算术够不着的地方。}瞬时速度是位移除以零；曲边区域的面积没有边长可乘；无穷多个数相加不是加法定义过的运算。它给出的办法不是硬算，而是先把``逼近''本身变成一个可检验的承诺，再在这个承诺之上重建除法与加法。

这一部分因此有两条并行的线。一条是内容线：极限、导数、积分、级数、多元。另一条是判断线，也是更该带走的那条——每个结论各自依赖什么条件，条件塌掉时会以什么方式出错。

\begin{center}
\begin{tikzpicture}[x=1cm,y=1cm,font=\small,>=stealth]
  \draw[black!45,fill=blue!8] (0,1.1) rectangle (3.4,2.3);
  \node[align=center] at (1.7,1.7) {整体量\\\(F\)};
  \draw[black!45,fill=blue!8] (0,-2.3) rectangle (3.4,-1.1);
  \node[align=center] at (1.7,-1.7) {局部率\\\(f=F'\)};
  \draw[->,black!65] (1.15,1.05) -- (1.15,-1.05);
  \node[left,black!75,align=center,font=\footnotesize] at (1.1,0) {求导\\看局部};
  \draw[->,black!65] (2.25,-1.05) -- (2.25,1.05);
  \node[right,black!75,align=center,font=\footnotesize] at (2.3,0) {积分\\做累积};
  \node[align=center,font=\footnotesize,black!60] at (1.7,-2.75) {两条箭头互逆 —— 微积分基本定理};
  \draw[->,dashed,black!50] (3.5,0) -- (5.3,0);
  \node[above,font=\footnotesize,black!65] at (4.4,0.1) {升到高维};
  \draw[black!45,fill=green!9] (5.4,-1.3) rectangle (11.6,1.3);
  \node[align=center] at (8.5,0.45) {区域内部的局部量累积起来，};
  \node[align=center] at (8.5,-0.15) {等于边界上的总效应};
  \node[align=center,font=\footnotesize,black!60] at (8.5,-0.85) {Green / Gauss / Stokes};
\end{tikzpicture}

\emph{左边这个来回是单变量微积分的全部骨架；右边只是同一句话在更高维度上的复述——边界定理并不是新原理，它是基本定理换了个形状。}
\end{center}

\subsection{五条贯穿始终的判断}

\textbf{极限不是``无限接近''，是一份可检验的承诺。}对手先报出容差 \(\varepsilon\)，你必须交出一个 \(\delta\) 兑现它。谁先出手、谁依赖谁，这个次序决定了一切：\(\delta\) 只能看 \(\varepsilon\)，就是一致连续；\(\delta\) 还能看点的位置，就只是逐点连续。后面遇到的收敛概念，包括级数收敛、函数列收敛、优化算法收敛，全都是在改写这份承诺的条款。含混地说``收敛了''，等于没说。

\textbf{导数是最佳线性近似的系数，不只是切线斜率。}把 \(f(a+h)=f(a)+f'(a)h+o(h)\) 当作定义来读，求导法则就不再是需要背的表格：它们是``把线性模型代进去、展开、扔掉高次项''的结果。链式法则说的是逐层放大率相乘——这句话就是反向传播的全部数学内容。到了多元情形，导数升级为线性映射，矩阵只是它在一组基下的写法。

\textbf{偏导数存在不等于可微。}这不是学究式的挑刺。可微意味着存在一个统一的线性近似，能同时对所有方向负责；而偏导数只沿坐标轴张望。存在一个函数，它在原点的每一个方向导数都存在，却依然不可微——因为方向导数没有线性地依赖方向。梯度之所以配叫``最速上升方向''，靠的正是可微这个前提。

\textbf{累积就是积分，而坐标一换，密度就得跟着改。}Riemann 和把``局部量乘小尺度''加起来，基本定理再把累积与变化率焊成互逆。换元法看起来是计算技巧，实质是坐标伸缩带来的密度因子——同一个因子在概率的变量替换、多重积分的 Jacobian 行列式、流模型的 \(\log|\det J|\) 里是同一件事。认出这一点，后面能省下大量重复学习。

\textbf{无限项的命运由尾部决定。}项趋于零远远不够，调和级数就是反例。各种判别法不是一张需要背的清单，它们是量同一段尾巴的几把不同尺子。这条判断的实用版本出现在随机逼近里：学习率序列要求 \(\sum\alpha_t=\infty\)（走得够远）同时 \(\sum\alpha_t^2<\infty\)（噪声可控），两个条件正好卡住尾部衰减的速度区间。

\subsection{六章怎么安排}

\hyperref[ch:011]{``极限与连续：不必到达，也能控制趋势''}是整部分的地基。第一次读 \(\varepsilon\)-\(\delta\) 时不必纠缠机械验算，先把``谁先给、谁后给''这个次序弄清楚，量词的形状自然就顺了。

\hyperref[ch:012]{``导数：用局部直线理解瞬时变化''}和\hyperref[ch:013]{``积分：把无数微小贡献累积起来''}是主线的两半，建议连读，因为基本定理只有在两边都熟悉之后才显出分量。

\hyperref[ch:010]{``函数与预备知识：先看清什么量依赖什么量''}更像随身工具，可以走快些；但陪域与值域的区别值得停一下，实践中不少事故就出在把``输出在 0 到 1 之间''当成了``取遍 0 到 1''。

\hyperref[ch:014]{``数列与级数：无限过程何时产生有限结果''}相对独立，急着往下走可以先跳过，等遇到收敛性讨论时回补。但 Taylor 余项那一篇不要跳——近似在多大范围内可信，是后面信赖域和二阶方法的直接依据。

\hyperref[ch:015]{``多元微积分：变化有方向，累积有几何''}是本部分与后续内容衔接最紧的一章。梯度、Jacobian、Hessian、Lagrange 乘子、变量替换，分别对应优化方向、反向传播、曲率与条件数、约束与正则化的对偶、以及概率密度变换。如果时间有限，这一章的前六篇比其余任何内容都更值得投入。

读的时候有个习惯值得保持：每遇到一条计算规则，回头问它依赖哪个极限、需要什么连续性或可微性、以及这些条件塌掉时会以什么方式出错。这一部分的例子刻意选得短，练的不是演算速度，是这个追问动作。
